% -*-memoria.tex-*-
% Este fichero es parte de la plantilla LaTeX para
% la realización de Proyectos Final de Carrera, protejido
% bajo los términos de la licencia GFDL.
% Para más información, la licencia completa viene incluida en el
% fichero fdl-1.3.tex

% Copyright (C) 2009 Pablo Recio Quijano 

%-------------------------------------------------------
% ---- Plantilla para libros / memorias PFC -----

% Realizada por Pablo Recio Quijano y Noelia Sales Montes 
% Formato de portada y primera página tomado del PFC de
% Francisco Javier Vázquez Púa, en su proyecto 'libgann'
% -------------------------------------------------------

\documentclass[a4paper,11pt]{book}

\usepackage{./estilos/estiloBase} % Basicamente son todas las
                                  % librerias usadas. En caso de que
                                  % falten librerias se van añadiendo
                                  % al fichero.
\usepackage{./estilos/colores}  % Algunos colores ya generados, para
                                % los algunos estilos más avanzados.
\usepackage{./estilos/comandos} % Algunos comandos personalizados

\usepackage{./estilos/dirtree}

\graphicspath{{./imagenes/}} % Indicamos la ruta donde se encuentran
                             % las imagenes, para ahorrarnos la ruta
                             % completa, y solo modificar aquí si en
                             % un momento dado lo movemos

\begin{document}

% Renombramos las figuras y las tablas
\renewcommand{\figurename}{Figura}
\renewcommand{\listfigurename}{Indice de figuras}
\renewcommand{\tablename}{Tabla}
\renewcommand{\listtablename}{Indice de tablas}

\pagestyle{empty}
% -*-portada.tex-*-
% Este fichero es parte de la plantilla LaTeX para
% la realización de Proyectos Final de Carrera, protejido
% bajo los términos de la licencia GFDL.
% Para más información, la licencia completa viene incluida en el
% fichero fdl-1.3.tex

% Fuente tomada del PFC 'libgann' de Javier Vázquez Púa

\begin{titlepage}

  \begin{center}

    \includegraphics[scale=0.35]{upv.png} \\
    
    \vspace{2.0cm}
    
    \LARGE{\textbf{E.T.S.I. TELECOMUNICACIÓN}} \\
    
    \vspace{1.0cm}
    
    \Large{\textbf{INGENIERÍA SUPERIOR EN TELECOMUNICACIÓN}} \\
    
    \vspace{3.0cm}
%mimateo: ¿seguro que quedamso en poner este título? 
%lluis: No, está esto puesto solamente de forma temporal. 
    \Large{ARDUINO Y MIPS} \\
    
    \vspace{2.0cm}
    
    \Large{Lluís Satorre González} \\
  
    \vspace{0.5cm}

    \large{\today}
    
  \end{center}
\end{titlepage}

\cleardoublepage

% -*-primerahoja.tex-*-
% Este fichero es parte de la plantilla LaTeX para
% la realización de Proyectos Final de Carrera, protejido
% bajo los términos de la licencia GFDL.
% Para más información, la licencia completa viene incluida en el
% fichero fdl-1.3.tex

% Fuente tomada del PFC 'libgann' de Javier Vázquez Púa

\begin{center}

  \includegraphics[scale=0.35]{upv.png} \\

  \vspace{2.0cm}

  \Large{ESCUELA SUPERIOR DE INGENIERÍA} \\

  \vspace{1.0cm}

  \large{INGENIERO SUPERIOR EN TELECOMUNICACIÓN} \\

  \vspace{2.0cm}

  \large{PLANTILLA PARA PROYECTO DE EJEMPLO} \\

  \vspace{1.0cm}

\end{center}

\begin{itemize}
\item \large{Departamento: Sistemas Informáticos y Computación}
\item \large{Director del proyecto: Miguel Ángel Mateo Pla}
\item \large{Autor del proyecto: Lluís Satorre González}
\end{itemize}

\vspace{1.0cm}

\begin{flushright}
  \large{Cádiz, \today} \\

  \vspace{2.5cm}

  \large{Fdo: Pablo Recio Quijano}
\end{flushright}

\cleardoublepage
\pagestyle{plain}

\frontmatter % Introducción, índices ...

% -*-previo.tex-*-
% Este fichero es parte de la plantilla LaTeX para
% la realización de Proyectos Final de Carrera, protejido
% bajo los términos de la licencia GFDL.
% Para más información, la licencia completa viene incluida en el
% fichero fdl-1.3.tex

% Copyright (C) 2009 Pablo Recio Quijano 

\section*{Agradecimientos}

Me gustaria agradecer y/o dedicar este texto a ...

\cleardoublepage

\section*{Licencia} % Por ejemplo GFDL, aunque puede ser cualquiera

Este documento ha sido liberado bajo Licencia GFDL 1.3 (GNU Free
Documentation License). Se incluyen los términos de la licencia en
inglés al final del mismo.\\

Copyright (c) 2009 Pablo Recio Quijano.\\

Permission is granted to copy, distribute and/or modify this document under the
terms of the GNU Free Documentation License, Version 1.3 or any later version
published by the Free Software Foundation; with no Invariant Sections, no
Front-Cover Texts, and no Back-Cover Texts. A copy of the license is included in
the section entitled "GNU Free Documentation License".\\

\cleardoublepage

\section*{Notación y formato}

Aquí incluiremos los aspectos relevantes a la notación y el formato a
lo largo del documento. Para simplificar podemos generar comandos
nuevos que nos ayuden a ello, ver \texttt{comandos.sty} para más
información. 

Cuando nos refiramos a un programa en concreto, utilizaremos la
notación: \\ \programa{emacs}.\\

Cuando nos refiramos a un comando, o función de un lenguaje, usaremos
la notación: \\ \comando{quicksort}.\\
\cleardoublepage

\tableofcontents
\listoffigures
\listoftables

\mainmatter % Contenido en si ...

\chapter{Motivación y contexto del proyecto}
% -*-cap1.tex-*-
% Este fichero es parte de la plantilla LaTeX para
% la realización de Proyectos Final de Carrera, protejido
% bajo los términos de la licencia GFDL.
% Para más información, la licencia completa viene incluida en el
% fichero fdl-1.3.tex

% Copyright (C) 2009 Pablo Recio Quijano 

\section{Contexto} % (fold)
% Arduino
% Microcontroladores
% Internet of Things
% Segunda revolución del garaje
%
\label{sec:Contexto}
Un microcontrolador es un circuito integrado programable que se compone de varios bloques funcionales, ente los que encontramos una CPU, memoria y periféricos. Con ayuda de un número mínimo de componentes externos, el microcontrolador permite crear pequeños circuitos que realicen tareas determinadas, que pueden ir desde la creación de sensores inteligentes hasta facilitar el manejo de un avión teledirigido. Un microprocesador, por contra, no incorpora memoria programable ni periféricos, por lo que necesita un mayor número de componentes externos para poder funcionar.

%mimateo: ¡¡Cuidado!!
%mimateo: El 4004 fue el primer microporcesador pero no el primer microcontrolador... Buscando en Google dan ese honor al TMS 1000 de Texas Instruments

En 1974 Gary W. Bone, en Texas Instruments, creó el primer microcontrolador de la histora: el TMS 1000 de 4 bits, después de varias disputas legales con Intel~\footnote{http://www.datamath.org/Story/Intel.htm}. Desde entonces, los microcontroladores han ido evolucionando y en la actualidad existen multitud de modelos diferentes entre los que elegir basados en diversas arquitecturas como pueden ser AVR, ARM-M o MIPS.\@

Un microcontrolador ejecuta las instrucciones del programa que se ha grabado en su memoria. Estos programas estaban elaborados originalmente en código ensamblador, aunque en la actualidad la mayoría se programa usando el lenguaje C. Para programar un microcontrolador hay que estar familiarizado con el funcionamiento del modelo concreto que se vaya a programar. No es lo mismo programar un microcontrolador ARM que un PIC32 o un Atmel AVR.\@ A pesar de que en todos ellos se utiliza C como principal lenguaje de programación, la velocidad de ejecución, las estructuras de traducción y los dispositivos internos difieren, en muchos casos sustancialmente. Por ejemplo, para escribir en un pin de un GPIO puede ser necesario configurar diferentes registros dependiendo del microcontrolador.

%mimateo: yo lo pondría al revá: primero el problema y luego la propuesta de solución
Para minimizar estos inconvenientes, en el año 2003, se comenzó a trabajar en la plataforma Wiring~\footnote{http://wiring.org.co/},  en el instituto IVREA, en Italia. Wiring proporcionaba una capa de abstracción sobre las funciones y registros de algunos microcontroladores. Esta plataforma está ideada con los diseñadores y artistas en mente ya que las plataformas de prototipado electrónico de aquel entonces estaban enfocadas a gente con una formación técnica en este ámbito. Con Wiring se pretendía reducir la dificultad del diseño electrónico y la programación. Dos años más tarde, nació Arduino que, junto con Wiring, ofrecía una plataforma de hardware y software abierto con la que cualquier persona, con unos conocimientos básicos de programación, podía crear sus propios proyectos.
% me falta una referencia en WIRING...

% section Contexto (end)

\section{Introducción}
Como alumno de la Escuela Técnica Superior de Ingenieros de Telecomunicación he tenido la oportunidad de aprender a programar diferentes dispositivos y arquitecturas en varias asignaturas. Comencé usando el lenguaje C en ordenadores con arquitectura x86 en la asignatura Programación del primer curso.  Más adelante, en el tercer curso en la asignatura `Sistemas Electrónicos Digitales', programé en ensamblador del microprocesador Motorola 68000.

\begin{figure}[H]
\begin{center}
\begin{tikzpicture}[
		decoration={
			markings,
			mark=at position 1 with {\arrow[ultra thick]{latex}},
		},
		path/.style={
			postaction=decorate,
		},
%		every node/.style={font=\sffamily}
	]
	% Boxes
	\node[
		draw,
		xshift=-6cm,
		minimum width=2cm,
	] (hid_in) {HID};
	\node[
		draw,
		minimum width=2cm,
		minimum height=2cm,
	] (pc) {PC};
	\node[
		draw,
		xshift=6cm,
		minimum width=2cm,
	] (hid_out) {HID};
	\node[
		draw,
		xshift=-6cm,
		yshift=1cm,
		minimum width=2cm,
] (mem_in) {Memoria};
	\node[
		draw,
		xshift=6cm,
		yshift=1cm,
		minimum width=2cm,
] (mem_out) {Memoria};
	\node[
		draw,
		xshift=-6cm,
		yshift=-1cm,
		minimum width=2cm,
] (net_in) {Red};
	\node[
		draw,
		xshift=6cm,
		yshift=-1cm,
		minimum width=2cm,
] (net_out) {Red};

% Conexiones
\draw[path] (mem_in.east) -- ++(2,0) node[above,pos=0.3]{digital} |- ([yshift=-0.5cm] pc.north west);
\draw[path] (hid_in.east) -- ++(2,0) node[above,pos=0.3]{digital} |- (pc.west);
\draw[path] (net_in.east) -- ++(2,0) node[below,pos=0.3]{digital} |- ([yshift=0.5cm] pc.south west);

\draw[path] ([yshift=-0.5cm] pc.north east) -- ++(2,0)  |- node[above,pos=0.7]{digital}(mem_out.west);
\draw[path] (pc.east) -- ++(2,0) |- node[above,pos=0.7]{digital} (hid_out.west);
\draw[path] ([yshift=0.5cm] pc.south east) -- ++(2,0) |- node[below,pos=0.7]{digital} (net_out.west);

\end{tikzpicture}
\end{center}
\caption{Origen y Destino de los datos procesados en un Ordenador}
\label{io_pc}
\end{figure}
%mimateo: lo de STDIN no estoy seguro ni que de naciera en UNIX. Lo que sí tengo claro es que se usa en muchos más sistemas operativos, de ahí la modificación que he realizado
Cuando se empieza a programar, normalmente se hace en un ordenador personal, con un sistema operativo de propósito general que ofrece el concepto de máquina extendida. En este entorno, el sistema operativo se encarga de gestionar los recursos disponibles en el sistema. De esta forma, cuando se escribe un programa, no es necesario acceder a bajo nivel a los dispositivos, tan solo hay que usar las bibliotecas que nos proporciona el sistema operativo y que usarán los drivers que alguien ha creado para usar el dispositivo en el sistema operativo, típicamente ese alguien es el propio fabricante del dispositivo.

Con el M68000 tuve la oportunidad de programar al más bajo nivel, al no disponer de un sistema operativo, accediendo directamente a los dispositivos a las rutinas de manejo de interrupciones. En lugar de obtener los datos desde un fichero o la entrada estándar (`STDIN'\footnote{Normalmente está asociado al teclado en un programa de terminal, teclado controlado por ese programa o por el sistema operativo.}) lo hacía directamente desde un teclado, desde el puerto serie o incluso a través de un convertidor analógico-digital (ADC), dispositivos que era necesario configurar ya que, debido a la falta de sistema operativo, no disponía de las funciones necesarias para controlarlos.
 
\begin{figure}[H]
%mimateo: en este figura falta una salida analógica pura
\begin{center}
\begin{tikzpicture}[
		decoration={
			markings,
			mark=at position 1 with {\arrow[ultra thick]{latex}},
		},
		path/.style={
			postaction=decorate,
		},
%		every node/.style={font=\sffamily}
	]
	% Boxes
	\node[
		draw,
		xshift=-6cm,
		minimum width=2cm,
	] (hid_in) {HID};
	\node[
		draw,
		xshift=-6cm,
		yshift=2.5cm,
		minimum width=2cm,
		text width = 3.5cm,
	] (sensores) {Sensores: audio, temperatura, humedad, luz \dots};
	\node[
		draw,
		xshift=-6cm,
		yshift=-2.5cm,
		minimum width=2cm,
		text width = 3.5cm,
	] (conexiones) {Conexiones Serie/paralelo: SPI, UART, Wire, CAN \dots};
	\node[
		draw,
		xshift=-0.5cm,
		minimum width=2cm,
		minimum height=3.5cm,
	] (uc) {\(\mu\)C};
	\node[
		draw,
		xshift=5cm,
		yshift=2.5cm,
		minimum width=2cm,
		text width=3.5cm,
	] (otros) {Otros elementos digitales: LEDs, Relés \dots};
	\node[
		draw,
		xshift=5cm,
		minimum width=2cm,
	] (hid_out) {HID};
	\node [
		draw,
		xshift=-6cm,
		yshift=1cm,
		minimum width=2cm,
	] (mem_in) {Memoria};
	\node [
		draw,
		xshift=5cm,
		yshift=1cm,
		minimum width=2cm,
	] (mem_out) {Memoria};
	\node[
		draw,
		xshift=-6cm,
		yshift=-1cm,
		minimum width=2cm,
	] (net_in) {Red};
	\node[
		draw,
		xshift=5cm,
		yshift=-1cm,
		minimum width=2cm,
	] (net_out) {Red};
	\node[
		draw,
		xshift=5cm,
		yshift=-2cm,
		minimum width=2cm,
	] (pwm) {PWM};
	\node[
		draw,
		xshift=4.5cm,
		yshift=-3.5cm,
		minimum width=2cm,
		text width=3.5cm,
	] (otros_analog) {Elementos analógicos: Altavoces, Buzzers, LEDs \dots};

% Conexiones
\draw[path] (sensores.east) -- ++(2,0) node[above,pos=0.5]{analógico} |- ([yshift=-0.5cm] uc.north west);
\draw[path] (mem_in.east) -- ++(2,0) node[above,pos=0.3]{digital} |- ([yshift=-1cm] uc.north west);
\draw[path] (hid_in.east) -- ++(2,0) node[above,pos=0.3]{digital} |- (uc.west);
\draw[path] (net_in.east) -- ++(2,0) node[below,pos=0.3]{digital} |- ([yshift=1cm] uc.south west);
\draw[path] (conexiones.east) -- ++(2,0) node[below,pos=0.5]{digital} |- ([yshift=0.5cm] uc.south west);

\draw[path] ([yshift=-0.5cm] uc.north east) -- ++(1,0)  |- node[above,pos=0.7]{digital}(otros.west);
\draw[path] ([yshift=-1cm] uc.north east) -- ++(2,0)  |- node[above,pos=0.7]{digital}(mem_out.west);
\draw[path] ([yshift=-1.5cm]uc.north east) -- ++(2,0) |- node[above,pos=0.7]{digital} (hid_out.west);
\draw[path] ([yshift=1.5cm] uc.south east) -- ++(2.5,0) |- node[below,pos=0.7]{digital} (net_out.west);
\draw[path] ([yshift=1cm] uc.south east) -- ++(2.25,0) |- node[below,pos=0.7]{digital} (pwm.west);
\draw[path] (pwm.east) -- ++(1,0) node[above]{analógico} |-  (otros_analog.east);
\draw[path] ([yshift=0.5cm] uc.south east) -- ++(1.75,0) node[below,pos=0.5]{analógico} |- (otros_analog.west);

\end{tikzpicture}
\end{center}
\caption{Origen y Destino de los datos procesados en un Microcontrolador.}
\label{io_uc}
\end{figure}
Esta plataforma tenía dos inconvenientes que hacían inviable su uso fuera de un entorno de aprendizaje o profesional.  Primero, era necesario programar todo en ensamblador, no pudiendo usar las abstracciones que nos otorga un lenguaje como C, o al menos no se podía hacer de una forma sencilla.  Además, el kit que se usaba en esta asignatura (SED) es difícil de encontrar y, en el caso de hacerlo, su precio es prohibitivo, superando los 1000\$\footnote{http://www.promaxprolink.com/training/tm683.htm}.  Ese mismo año también realicé prácticas de programación con VHDL en FPGAs que también tenían los mismos problemas que el Motorola 68000: alto precio y programación de muy bajo nivel.

% sería interesante que dijeras el año exacto o la edición de las Google IO a las que te refieres
Un año más tarde conocí la plataforma Arduino, a raíz del lanzamiento del Arduino ADK durante la Google I/O del 2011.\@ La plataforma presentada se componía de una placa con un microcontrolador Atmel2560 de 8 bits y era de un tamaño mucho menor que los kits que había usado a lo largo de la carrera.  El precio era de 60~\euro{} para ese modelo, que era el más caro en ese momento. Aún así, era mucho más asequible que las plataformas de prototipado que había utilizado hasta ese momento.  

Uno de los aspectos más importantes de esta plataforma es que no es necesario conocer el microcontrolador para llevar a cabo proyectos en Arduino.  Por ejemplo, para programar el parpadeo de un LED todo lo que hay que averiguar es a qué pin (numerados en la placa) está conectado el diodo LED.\@ Se puede programar usando C o C++ y cuenta con un entorno de desarrollo que permite la utilización de múltiples funciones que hacen posible programar a un nivel más alto que al que se haría si usáramos las herramientas que proporciona Atmel para sus microcontroladores.  

El hardware y el software de Arduino son abiertos.  Es posible obtener los esquemáticos de cualquier placa Arduino desde su página web y la biblioteca principal (en la que se basa cualquier proyecto Arduino) viene incluida con el entorno de programación como código fuente.  Esto significa que existe la posibilidad de comprar una placa oficial a través de la tienda de Arduino o algún distribuidor, pero también es posible construirla uno mismo o encontrar placas personalizadas a través de Internet que lleven caracterísiticas no incluidas en las oficiales o que sean más baratas.

A esto se suma la posibilidad de utilizar el entorno Arduino y muchas de sus bibliotecas de funciones incluso con placas que no estén basadas en los mismos procesadores que las oficiales. En este tipo de placas se debe distribuir algún compilador específico para las mismas así como una implementación de las bibliotecas básicas de funcionamiento de Arduino.

Un ejemplo de este tipo de placas es la que he empleado en este trabajo: ChipKIT MAX32.  A diferencia de las placas Arduino convencionales, la gama ChipKIT de Digilent tiene como microcontrolador un PIC de 32 bits en vez de un Atmega de 8 bits.  Está placa ofrece más potencia y más prestaciones como, por ejemplo, la posibilidad de depurar nuestro programa gracias al conector JTAG que lleva incorporado. Las diferencias están comentadas más a fondo en el siguiente capítulo.

\section{Objetivos}
La utilización de una placa de la gama chipKit implica la utilización de una nueva herramienta de programación, mucho más potente que la original de Arduino. En esta plataforma disponemos de un control casi completo de la forma en la que se generarán los archivos compilados y es en este tipo de control en los que hemos centrado este PFC.\@

Teniendo en cuenta esta orientación al estudio a bajo nivel de los programas gnerados con el entorno Arduino, podemos dividir los objetivos de la siguiente forma:

\subsubsection{Analizar el proceso de compilación de \programa{MPIDE}}
Lo primero que realizaré será un análisis del proceso que se sigue en el entorno de desarrollo de ChipKIT desde que escribimos el programa hasta que este esté funcionando en la placa microcontroladora.  Esto incluye todas las órdenes o programas que intervengan en este proceso y sus respectivas opciones, además de los ficheros involucrados.

Para lograr flexibilidad en este proceso, propondré una alternativa al IDE, que hará uso de Makefiles y me permitirá realizar todo el proceso anterior con una sola orden. También podré generar resultados intermedios como, por ejemplo, el código compilado, pero no ensamblado, y podré modificar las opciones de cada una de las órdenes.

\subsubsection{Comparar el efecto de los distintos niveles de optimización}
El compilador ofrece diferentes niveles de optimización y según el nivel se llevan a cabo unas u otras optimizaciones. Uno de los objetivos será identificar estas optimizaciones y ver su efecto en el código. Usando una aplicación desarrollada para el microcontrolador pondré de manifiesto las diferencias entre los distintos niveles de optimización, qué ventajas y qué desventajas tienen cada uno de ellos.

\subsubsection{Poner de manifiesto las limitaciones de la biblioteca Arduino}
La biblioteca Arduino nos ofrece una capa de abstracción sobre las funciones específicas del microcontrolador pero que tiene limitaciones. Comprobaré la diferencia entre usar esta biblioteca o acudir a las bibliotecas del fabricante; en este caso, Microchip.

\subsubsection{Desarrollar una aplicación enfocada en el ámbito de las telecomunicaciones}
Desarrollaré una aplicación en la que se pondrán en práctica conceptos relacionados con las telecomunicaciones, aprendidos durante la carrera. La función de la aplicación será detectar tonos mediante el uso del algoritmo de Goertzel.


\section{Estructura de la memoria}
Este trabajo se organiza en diferentes capítulos en los que trataré temas específicos.

En el capítulo~\ref{cap2} hablaré sobre los dispositivos hardware empleados en el trabajo. Esto incluirá una descripción de la placa microcontroladora elegida junto con los accesorios que la acompañan. La descripción de la placa incluirá también un apartado dedicado a su chip microcontrolador.

En el capítulo~\ref{cap3} describo el apartado software de la plataforma Chipkit. Realizaré una descripción de los programas que se pueden usar para programar una placa Chipkit: \programa{MPIDE} y \programa{MPLAB}. En este capítulo también trataré el primero de los objetivos, ya que analizaré el proceso de compilación de \programa{MPIDE} y propondré una alternativa al mismo.

Usando los elementos tratados en los dos capítulos anteriores, en el capítulo~\ref{cap4} se llevan a cabo el resto de objetivos planteados. Comenzaré tratando los niveles de optimización y sus distintos efectos. Pondré a prueba los distintos niveles con la ayuda de una aplicación que corra en el Chipkit Max32. Una vez hecho esto comentaré las posibles limitaciones de la plataforma Arduino y como es posible solventar alguna de ellas. Para finalizar el capítulo describiré una aplicación orientada a la carrera. Como ya he comentado en los objetivos, se tratará de la implementación de un detector de tonos usando el algoritmo de Goertzel, una aplicación típica de DSPs pero que se puede implementar también en el Max32.

Por último, en el capítulo~\ref{cap5} comentaré mis conclusiones, tanto técnicas como personales, sobre lo que ha supuesto este trabajo.

Al final de la memoria también se pueden encontrar varios apéndices que desarrollan más detalladamente ciertos temas como, por ejemplo, la descripción detallada del archivo Makefile que utilizo como alternativa al IDE de Chipkit o el funcionamiento del programa `Sorted' que uso en el capítulo~\ref{cap4} para medir las diferencias entre niveles de optimización.


\chapter{Desarrollo del proyecto}
% -*-cap2.tex-*-
% Este fichero es parte de la plantilla LaTeX para
% la realización de Proyectos Final de Carrera, protejido
% bajo los términos de la licencia GFDL.
% Para más información, la licencia completa viene incluida en el
% fichero fdl-1.3.tex

% Copyright (C) 2009 Pablo Recio Quijano 

En este capitulo analizaré las características del Arduino Mega 2560\footnote{En la introducción menciono el Arduino Mega ADK, pero la única diferencia entre ambos es que este último puede hacer de Host USB para poder conectar un teléfono móvil Android.} y del ChipKIT MAX32, que es el equivalente en chipKIT, tanto en formato y dimensiones como en precio.

\section{Arduino Mega 2560}
La familia Arduino se compone de multitud de placas microcontroladoras con características diferentes pero hay dos que se pueden destacar sobre las demas debido a que son las más utilizadas.  Estas dos placas son la Arduino Uno y Arduino Mega 2560. La placa Arduino Uno es la más básica y económica, con un formato reducido, mientras que la segunda tiene un precio superior y mejores características.

\figura{ArduinoMega2560_R3_Front.jpg}{scale=0.1}{Arduino Mega 2560}{mega2560}{H}

Arduino Mega 2560 es una placa microcontroladora basada en el chip ATmega2560 que fue introducida en el mercado durante el último trimestre de 2010.  
Es una actualización del Arduino Mega, basada en el chip ATmega1280.\@ Dispone de 54 pines de entrada/salida (14 de estos pines pueden ser usado como salidas PWM), 16 entradas analógicas, 4 UARTs (puertos serie que se pueden usar para conectar la placa al ordenador, a periféricos o a otra placa microcontroladora), un oscilador de cristal de 16 MHz, un conector de alimentación, un botón de reinicio y un conector ICSP.\@

\begin{table}
\begin{center}
\begin{tabular}{| l | c |}
\hline
Microcontrolador 						& ATmega2560 \\
Voltaje									& 5~V \\
Voltaje de alimentación (recomendado) 	& 7-12~V \\
Voltaje de alimentación (límites) 		& 6-20~V \\
Pines Entrada/Salida digitales			& 54 (15 salidas PWM) \\
Pines de entrada analógicos				& 16 \\
Corriente por cada pin entrada/salida 	& 40~mA \\
Corriente en el pin de 3,3~V				& 50~mA \\
Memoria Flash 							& 256~kB (8 kB usados por el bootloader) \\
SRAM									& 8~kB \\
EEPROM									& 4~kB \\
Frecuencia del reloj					& 16~MHz \\
\hline
\end{tabular}
\end{center}
\caption{Especificaciones de Arduino Mega 2560.}
\label{tab:mega2560_specs}
\end{table}
% Cosas que mencionar
% Cantidad de memoria
% Voltaje (5V)
% JTAG Pero en las lineas de los ADC
% ADC (pag 275 datasheet) 'solo' 15ksps
Una de las carácterísticas de las placas Arduino es su voltaje de funcionamiento de 5~V. Este voltaje se puede obtener desde la conexión USB al PC ya que la línea de alimentación de este es también de 5~V, o se puede usar una fuente de alimentación externa con un voltaje de entre 7~V y 12~V.

El Arduino Mega 2560 dispone de 256~kB de memoria flash.\footnote{En los microcontroladores de Atmel esto viene indicado en el propio nombre, en este caso 2560 hace referencia a sus 256~kB de memoria Flash.} Esta cantidad es suficiente para la mayoría de proyectos a los que estaría destinada esta placa. Por otro lado su memoria SRAM es de tan solo 8~kB, lo que puede ser el factor limitante en algunos proyectos.

Esta placa tiene también la posibilidad de depurar sus programas a través del puerto JTAG del Atmega 2560. En lugar de tener una conexión dedicada, los pines del puerto JTAG están situados en los mismos dedicados a las entradas analógicas, lo que dificulta la depuración de código que haga uso de estas entradas.

Otra característica que se puede destacar de este Arduino son los convertidores analógico-digital que tiene incorporados el microcontrolador. Hay 2 ADCs en los que podemos encontrar 8 pines de entrada multiplexados en cada uno. Su resolución es de 10~bits y su máxima frecuencia de muestreo (con toda la resolución) es de 15~ksps. Como veremos después, este valor es bajo y puede resultar determinante dependiendo del proyecto en que se vaya a usar. Es posible obtener una mayor frecuencia de muestreo, pero a cambio de reducir la resolución.

\subsection*{ Atmel Atmega2560}
La placa Mega 2560 recibe su nombre del microcontrolador que lleva instalado, un Atmel Atmega2560. Todos los periféricos que están incluidos en la placa los podemos encontrar integrados en este chip, de forma que el resto de componentes electrónicos son los dedicados a las entradas y salidas o a la alimentación de la propia placa, con la excepción del chip Atmega 16u2 que se encarga de gestionar la conexión a través del puerto USB con el ordenador.

El núcleo del Atmega 2560 está basado en una arquitectura AVR de 8~bits. AVR utiliza una arquitectura Harvard, de forma que la CPU dispone de buses independientes para datos y programa. El ciclo de instrucción está segmentado en dos partes: mientras una instrucción es ejecutada, la siguiente está siendo obtenida desde la memoria Flash. De esta forma, se ejecuta una instrucción por ciclo, siempre y cuando no haya dependencias entre las instrucciones.
% Cosas para mencionar
% Frecuencia de reloj
% 8 Bits
% Arquitectura AVR
% Pipeline de 2 fases
% NO MDU
% NO FPU
% 32 Registros
% No MMU
\figura{atmel2560_diagrama.jpg}{scale=0.7}{Diagrama de la arquitectura AVR}{atmel2560}{H}
Como se puede ver en la figura~\ref{atmel2560}, el núncleo AVR consta de 32 registros de 8 bits de propósito general en los que se almacenan los datos que serán procesados en la unidad aritmético-lógica (ALU). Esta unidad soporta operaciones aritméticas y lógicas entre registros o entre un registro y una constante guardando el resultado de nuevo en los registros, en un solo ciclo. Estas operaciones son llevadas a cabo sobre operandos de 8 bits (el tamaño de cada registro) por lo que cualquier operación sobre números enteros de 32~bits (operación más habitual en programas) requerirá más ciclos. Las operaciones aritméticas que puede realizar esta ALU son suma, resta y multiplicación de números de coma fija; sin embargo no dispone de una instrucción de división, por lo que es necesario usar una que funcione mediante software. Tampoco tiene una unidad de coma flotante (FPU) por lo que estamos limitados a números en coma fija o a la emulación de coma flotante a través de bibliotecas, que es lo que sucede cuando usamos variables de tipo `float' durante el desarrollo de una aplicación para Arduino.

Este microcontrolador incluye varios perífericos. Posee un total de 6 temporizadores/contadores, dos de ellos de 8~bits y los otros 4 de 16~bits de resolución. También tiene un RTC (Real Time Counter) que permite llevar la cuenta del tiempo real con un oscilador externo conectado al pin TOSC1. Si observamos el esquemático del Arduino Mega 2560 vemos que ese pin no está conectado a nada, es decir, no contiene un oscilador, y tampoco está asociado con ninguno de los conectores disponibles en la placa, por lo que no podremos usar esta funcionalidad. Otros periféricos que encontramos en este chip son 2 interfaces serie: SPI y 2-wire, y un temporizador programable `watchdog' que requiere también un oscilador externo, por lo que tampoco podremos hacer uso de esta característica con la placa Arduino. Por último, también tiene la capacidad de generar interrupciones cuando algún pin cambia su valor, por lo que no será necesario estar activamente esperando durante la ejecución del programa.

\section{ChipKIT Max32}
La placa microcontroladora elegida para el trabajo es una ChipKIT Max32 de DigilentINC.\@ Es el modelo más avanzado de la familia ChipKIT, que se compone además del Uno32, uC32 y DP32.

\figura{chipKIT-Max32.jpg}{scale=0.6}{Chipkit Max32}{max32}{H}

En el centro del Max32 encontramos como microcontrolador un Microchip PIC32MX795F12L que funciona a 80~MHz y con el que disponemos de 512~KB de memoria flash y 128~KB de memoria RAM.\@ Sin embargo, a diferencia del Arduino Mega 2560, no incluye memoria EEPROM en el chip por lo que no podemos almacenar datos de forma no volátil. 

El Max32 tiene 83 puertos de entrada/salida además de otros dispositivos como UART, SPI e I2C. Dieciséis de los 83 puertos pueden ser usados como entradas analógicas o como entradas/salidas digitales. A diferencia de su equivalente Arduino (Mega), el ChipKIT Max32 también tiene integrado en su MCU 1 módulo 10/100 Ethernet, 1 controlador USB 2.0 y 2 controladores CAN, aunque para hacer uso de estos dispositivos necesitaremos el ChipKIT Network Shield ya que es en esta placa donde están los conectores. 

Al igual que el Arduino Mega 2560, el microcontrolador de esta placa también tiene un RTC (Real Time Clock) que nos permita mantener la hora actual siempre que esté conectado a alguna fuente de alimentación. Para poder usar esta característica es necesario conectar un cristal de cuarzo que funcione a 32~768~kHz al pin correspondiente. En este caso, a diferencia del Arduino Mega 2560, el pin del oscilador se encuentra mapeado a uno de los pines de entrada/salida en los conectores de la placa por lo que podremos hacer uso de este periférico.\footnote{En el ChipKIT Uno32 podemos encontrar dos agujeros en la placa donde es posible soldar directamente un cristal para el RTC.}

\begin{table}[H]
\begin{center}
\begin{tabular}{| l | c |}
\hline
Microcontrolador 						& Microchip PIC32MX795F512L \\
Voltaje									& 3,3~V \\
Voltaje de alimentación (recomendado) 	& 7-15~V \\
Voltaje de alimentación (límites) 		& 6-20~V \\
Pines Entrada/Salida digitales			& 83 (14 salidas PWM) \\
Pines de entrada analógicos				& 16 \\
Corriente por cada pin entrada/salida 	& 18 mA \\
Corriente en el pin de 3,3~V				& N/A  \\
Memoria Flash 							& 512~kB (8 kB usados por el bootloader) \\
SRAM									& 128~kB \\
EEPROM									& No \\
Frecuencia del reloj					& 80~MHz \\
\hline
\end{tabular}
\end{center}
\caption{Especificaciones de ChipKIT MAX32.}
\label{tab:max32_specs}
\end{table}

Este dispositivo, al igual que el resto de microcontroladores de la familia ChipKIT, funciona con un voltaje lógico de 3,3~V frente a los microcontroladores Arduino como el Uno o el Mega que usan un voltaje de 5~V. Es importante tener este dato en cuenta a la hora de usar shields o dispositivos pensados para placas Arduino ya que están diseñandos para 5~V y es posible que no reciban bien los datos desde el Max32 por no superar en alguna ocasión el umbral de detección lógico, es decir, el voltaje a partir del cual se determina que el bit recibido es un `1'. Por otro lado esta placa puede recibir señales con voltaje de 5~V sin dañarse por lo que no supone ningún peligro conectar shields pensados para placas Arduino.

Podemos encontrar más información, además de esquemáticos y manuales en la página de información del ChipKIT Max32 en la web de Digilent.\cite{website:max32}

\subsection*{Microchip PIC32MX795F512L}
Este microcontrolador es el modelo más avanzado de la familia PIC32MX de Microchip. Como su nombre indica, se trata de un procesador de 32~bits, frente a los 8 del Atmel. Su arquitectura, aunque también basada en Harvard, es completamente diferente de la del Atmega. Se trata de la arquitectura MIPS, en su versión M4K. Este núcleo consta de 32 registros de propósito general de 32~bits. También cuenta con otro juego de 32 registros (shadow registers) que usa durante interrupciones. De esta forma, cuando se produce alguna interrupción no es necesario guardar los registros principales por lo que la transición de un estado a otro es más rápida.

\figura{m4k_diagram.jpg}{scale=0.7}{Diagrama del Núcleo M4K}{m4kdiagrama}{H}

Como se puede apreciar en la figura~\ref{m4kdiagrama} este chip cuenta con una MDU (unidad de multiplicación y división), un coprocesador y una unidad de administración de memoria (MMU). Ninguno de estas unidades podemos encontrarlas en el chip de Atmel. El MDU nos permite realizar multiplicaciónes entre operandos de 16~bits (o uno de 16 y otro de 32~bits) en un solo ciclo y entre dos operandos de 32~bits en dos o tres ciclos frente al ciclo que tarda la multiplicación de dos operandos de 8 bits en el núcleo AVR.\@ Además, con esta unidad se pueden realizar divisiones mediante hardware sin tener usar bibliotecas (como pasa con el Atmega), y el tiempo de estas operaciones es de 35 ciclos de reloj en el peor de los casos.

El coprocesador es una interfaz de registros que sirve para comunicar información de control y estado entre el programa y la CPU.\@ A través del coprocesador tenemos acceso al `core timer', al control de interrupciones y excepciones, a la configuración de la memoria virtual, al control de los registros `shadow', a la identificación del procesador y al control de depuración.

\figura{m4k_pipeline.jpg}{scale=0.7}{Segmentación (simplificada) en MIPS}{mipspipeline}{H}

MIPS fue desarrollada por John L. Hennessy y su equipo como una arquitectura de tipo RISC de 32~bits. Al igual que AVR, MIPS cuenta con segmentación de instrucciones, pero dividido en 5 etapas mostradas en la figura~\ref{mipspipeline}. Estas etapas son: Instruction Fetch, Execution, Memory Fetch, Align y Writeback. Podemos encontrar una descripción detallada de los procesos que se llevan a cabo en cada una de ellas en el apartado 2.3 del manual de referencia de PIC32MX accesible a través de la página web del PIC32MX795F512L~\cite{website:pic32mx}. Este microcontrolador también incluye mecanismos para evitar las interrupciones en la ejecución de instrucciones por dependencia de datos. Para ello utiliza realimentación de forma que la etapa de ejecución puede obtener operandos directamente desde la etapa `M' o la etapa `A' como se puede ver en la figura~\ref{mipspipeline}. Usando la realimentación se puede enviar los operandos desde una de esas etapas reduciendo posibles ciclos de espera. El incremento en rendimiento en operaciones que involucren a la ALU supone alcanzar una instrucción por ciclo incluso con la presencia de dependencia de registros.

Desde el punto de vista académico, usar este microcontrolador con arquitectura MIPS frente al Atmel y su AVR tiene la ventaja de que MIPS se estudia en la asignatura `Arquitectura de computadores y sistemas operativos I'. Esto ayuda a la hora de desarrollar programas al estar ya familiarizado con el juego de instrucciones y su segmentación de instrucciones, incluyendo los riesgos de control y datos, y cómo se reducen.

\section{Accesorios}
Durante el trabajo utilizaré varios accesorios junto con el ChipKIT MAX32 que me permiten depurar programas, conectar la placa a la red Ethernet y me proporcionan algunos elementos de entrada-salida como LEDs, botones o una pantalla OLED.\@
\subsection{Chipkit PGM}
Para la depuración de los programas con MPLAB usaré el programador chipKIT PGM\cite{website:pgm}, con el que podré usar el puerto JTAG del Max32. Este programador está diseñado para ser usado exclusivamente con la familia de dispositivos chipKIT, a diferencia de picKIT3~\cite{website:pickit} que puede ser usado con todos los microcontroladores de Microchip.

\figura{chipKIT-PGM.png}{scale=0.3}{ChipKIT PGM}{pgm}{H}

\subsection{ChipKIT Network Shield}

El ChipKIT Network Shield es, como su nombre indica, un shield para el ChipKIT MAX32.\@ Cabe destacar que este shield no sirve para el ChipKIT Uno32 puesto que el microcontrolador que está montado en ella no dispone de las características de las que hace uso esta placa, como Ethernet y host USB.\@ Se coloca encima del microcontrolador tal y como lo hace cualquier shield de Arduino.

\figura{chipKIT-NetworkShield.jpg}{scale=0.5}{Chipkit Network Shield}{networkshield}{H}

Este shield nos proporciona las conexiones necesarias para poder hacer uso del controlador Ethernet y USB que lleva el Max32 integrado y a los buses CAN.\@ Además de esto también tiene integrado un oscilador de 32~768 KHz que nos permite hacer uso del RTC que también tiene integrado el Max32 y una memoria EEPROM de 256~Kb a la que podemos acceder mediante I2C. En la página web del Network Shield~\cite{website:network_shield} podemos encontrar las bibliotecas que nos sirven para hacer uso de sus características, en el entorno ChipKIT.\@

Características del ChipKIT Network Shield:
\begin{itemize}
	\item Puede ser usado con el Max32
	\item SMSC LAN8720 10/100 Ethernet PHY
	\item Conector RJ45
	\item Dispositivo USB y conectores Host
	\item 2 transceptores CAN MCP2551
	\item 2 conectores de 12 pines para los buses CAN
	\item 2 conectores para cadena margarita (daisy chain) I2C
	\item EEPROM de 256~Kbits con interfaz I2~C
	\item Oscilador a 32~768~KHz
\end{itemize}

\subsection{ChipKIT Basic I/O Shield}
El otro shield que usaré será el ChipKIT Basic I/O Shield que nos servirá para interactuar físicamente con el microcontrolador ya que este shield nos proporciona diferentes métodos de entrada y salida de información.

\figura{chipKIT-BasicIOShield.jpg}{scale=0.5}{ChipKIT Basic I/O Shield}{ioshield}{H}

El ChipKIT Basic I/O shield dispone de los siguientes dispositivos:
\begin{itemize}
	\item EEPROM de 256~Kbits con interfaz I2C
	\item Pantalla OLED de 128$\times$32 píxeles
	\item Sensor de temperatura con interfaz I2C
	\item Conector para cadena margarita (daisy chain) I2C
	\item 4 drivers FET de colector abierto
	\item Potenciometro analógico
	\item 4 pulsadores y cuatro interuptores
	\item 8 LEDs
\end{itemize}

Este shield está inicialmente pensado para ser usado en un chipKIT Uno32, pero podemos conectarlo al Max32 sin ningún problema ya que comparten el mismo esquema de pines (aunque el Max32 tiene 83 pines frente a los 42 del Uno32). También es posible conectarlo al Network Shield y este a su vez al Max32, teniendo en cuenta la posibilidad de que algunas funciones en los dos shields puedan compartir pines.

Igual que el Network Shield, el Basic I/O Shield también dispone de una página web~\cite{website:io_shield} en la que podemos encontrar su manual, esquemáticos y las bibliotecas necesarias para usar en el entorno ChipKIT.\@


\backmatter % Apéndices, bibliografia ...


\clearpage
\addcontentsline{toc}{chapter}{Bibliografia y referencias}
\bibliographystyle{plain}
\bibliography{bibliografia}


\addcontentsline{toc}{chapter}{Software usado}
\chapter*{Software utilizado}
% -*-programas.tex-*-
% Este fichero es parte de la plantilla LaTeX para
% la realización de Proyectos Final de Carrera, protejido
% bajo los términos de la licencia GFDL.
% Para más información, la licencia completa viene incluida en el
% fichero fdl-1.3.tex

% Copyright (C) 2009 Pablo Recio Quijano 
Para la realización de esta memoria he utilizado diferentes programas.

\section*{Vim}
\programa{Vim} es uno de los editores de texto más utilizado en entornos Unix, enfocado en programación. La ventaja de usar este editor es que está preinstalado en cualquier sistema Unix (puede estar instalado Vim o Vi) y es altamente personalizable mediante su archivo de configuración y con plugins.

\section*{Sublime Text + Latexing}
Otro editor de textos que he utilizado para la creación de esta memoria es \programa{Sublime Text} que, a diferencia de \programa{Vim}, es de pago. Al igual que \programa{Vim} es posible ampliar su funcionalidad con plugins.

\section*{Git}
Para mantener un control de versiones he optado por utilizar \programa{Git}. Con \programa{Git} puedo, no solo guardar las versiones, sino también crear diferentes ramas o utilizar un servicio como Github.com, que me permite almacenar tanto esta memoria como los programas que he creado para este trabajo.

\section*{GNU Make}

\programa{GNU Make} es el programa de recompilación y de control de
dependencias por excelencia. Se puede utilizar para compilar proyectos
software en diversos códigos, o como en el caso de este documento,
para compilar documentos \LaTeX{} con diversas opciones.\\

Para más información \cite{pdf:make}


\addcontentsline{toc}{chapter}{Apéndice 1: Resultado de la compilación en \programa{MPIDE}}
\appendix
\chapter*{Apéndice 1: Resultado de la compilación en \programa{MPIDE}}\label{chap:apendice1}
\begin{lstlisting}[language=bash, label=bash:mpidecompilado, basicstyle=\small, breaklines=true]

/Applications/mpide.app/Contents/Resources/Java/hardware/pic32/compiler/pic32-tools/bin/pic32-g++  -O2  -c  -mno-smart-io  -w  -fno-exceptions  -ffunction-sections  -fdata-sections  -g  -mdebugger  -Wcast-align  -fno-short-double  -mprocessor=32MX795F512L  -DF_CPU=80000000L  -DARDUINO=23  -D_BOARD_MEGA_    -DMPIDEVER=0x01000305  -DMPIDE=23      -I/Applications/Mpide.app/Contents/Resources/Java/examples/1.Basics/Blink   -I/Applications/Mpide.app/Contents/Resources/Java/hardware/pic32/cores/pic32   -I/Applications/Mpide.app/Contents/Resources/Java/hardware/pic32/variants/Max32    /var/folders/f0/q5xc0_tj4jv4ljd_lfz4pg5h0000gn/T/build2117742867436922308.tmp/Blink.cpp  -o  /var/folders/f0/q5xc0_tj4jv4ljd_lfz4pg5h0000gn/T/build2117742867436922308.tmp/Blink.cpp.o
/Applications/mpide.app/Contents/Resources/Java/hardware/pic32/compiler/pic32-tools/bin/pic32-g++  -O2  -g1  -c  -mprocessor=32MX795F512L  -DF_CPU=80000000L  -DARDUINO=23  -D_BOARD_MEGA_    -DMPIDEVER=0x01000305  -DMPIDE=23      -I/Applications/Mpide.app/Contents/Resources/Java/hardware/pic32/cores/pic32   -I/Applications/Mpide.app/Contents/Resources/Java/hardware/pic32/variants/Max32    /Applications/Mpide.app/Contents/Resources/Java/hardware/pic32/cores/pic32/cpp-startup.S  -o  /var/folders/f0/q5xc0_tj4jv4ljd_lfz4pg5h0000gn/T/build2117742867436922308.tmp/cpp-startup.S.o
/Applications/mpide.app/Contents/Resources/Java/hardware/pic32/compiler/pic32-tools/bin/pic32-g++  -O2  -g1  -c  -mprocessor=32MX795F512L  -DF_CPU=80000000L  -DARDUINO=23  -D_BOARD_MEGA_    -DMPIDEVER=0x01000305  -DMPIDE=23      -I/Applications/Mpide.app/Contents/Resources/Java/hardware/pic32/cores/pic32   -I/Applications/Mpide.app/Contents/Resources/Java/hardware/pic32/variants/Max32    /Applications/Mpide.app/Contents/Resources/Java/hardware/pic32/cores/pic32/crti.S  -o  /var/folders/f0/q5xc0_tj4jv4ljd_lfz4pg5h0000gn/T/build2117742867436922308.tmp/crti.S.o
/Applications/mpide.app/Contents/Resources/Java/hardware/pic32/compiler/pic32-tools/bin/pic32-g++  -O2  -g1  -c  -mprocessor=32MX795F512L  -DF_CPU=80000000L  -DARDUINO=23  -D_BOARD_MEGA_    -DMPIDEVER=0x01000305  -DMPIDE=23      -I/Applications/Mpide.app/Contents/Resources/Java/hardware/pic32/cores/pic32   -I/Applications/Mpide.app/Contents/Resources/Java/hardware/pic32/variants/Max32    /Applications/Mpide.app/Contents/Resources/Java/hardware/pic32/cores/pic32/crtn.S  -o  /var/folders/f0/q5xc0_tj4jv4ljd_lfz4pg5h0000gn/T/build2117742867436922308.tmp/crtn.S.o
/Applications/mpide.app/Contents/Resources/Java/hardware/pic32/compiler/pic32-tools/bin/pic32-gcc  -O2  -c  -mno-smart-io  -w  -ffunction-sections  -fdata-sections  -g  -mdebugger  -Wcast-align  -fno-short-double  -mprocessor=32MX795F512L  -DF_CPU=80000000L  -DARDUINO=23  -D_BOARD_MEGA_    -DMPIDEVER=0x01000305  -DMPIDE=23      -I/Applications/Mpide.app/Contents/Resources/Java/hardware/pic32/cores/pic32   -I/Applications/Mpide.app/Contents/Resources/Java/hardware/pic32/variants/Max32    /Applications/Mpide.app/Contents/Resources/Java/hardware/pic32/cores/pic32/exceptions.c  -o  /var/folders/f0/q5xc0_tj4jv4ljd_lfz4pg5h0000gn/T/build2117742867436922308.tmp/exceptions.c.o
/Applications/mpide.app/Contents/Resources/Java/hardware/pic32/compiler/pic32-tools/bin/pic32-gcc  -O2  -c  -mno-smart-io  -w  -ffunction-sections  -fdata-sections  -g  -mdebugger  -Wcast-align  -fno-short-double  -mprocessor=32MX795F512L  -DF_CPU=80000000L  -DARDUINO=23  -D_BOARD_MEGA_    -DMPIDEVER=0x01000305  -DMPIDE=23      -I/Applications/Mpide.app/Contents/Resources/Java/hardware/pic32/cores/pic32   -I/Applications/Mpide.app/Contents/Resources/Java/hardware/pic32/variants/Max32    /Applications/Mpide.app/Contents/Resources/Java/hardware/pic32/cores/pic32/HardwareSerial_cdcacm.c  -o  /var/folders/f0/q5xc0_tj4jv4ljd_lfz4pg5h0000gn/T/build2117742867436922308.tmp/HardwareSerial_cdcacm.c.o
/Applications/mpide.app/Contents/Resources/Java/hardware/pic32/compiler/pic32-tools/bin/pic32-gcc  -O2  -c  -mno-smart-io  -w  -ffunction-sections  -fdata-sections  -g  -mdebugger  -Wcast-align  -fno-short-double  -mprocessor=32MX795F512L  -DF_CPU=80000000L  -DARDUINO=23  -D_BOARD_MEGA_    -DMPIDEVER=0x01000305  -DMPIDE=23      -I/Applications/Mpide.app/Contents/Resources/Java/hardware/pic32/cores/pic32   -I/Applications/Mpide.app/Contents/Resources/Java/hardware/pic32/variants/Max32    /Applications/Mpide.app/Contents/Resources/Java/hardware/pic32/cores/pic32/HardwareSerial_usb.c  -o  /var/folders/f0/q5xc0_tj4jv4ljd_lfz4pg5h0000gn/T/build2117742867436922308.tmp/HardwareSerial_usb.c.o
/Applications/mpide.app/Contents/Resources/Java/hardware/pic32/compiler/pic32-tools/bin/pic32-gcc  -O2  -c  -mno-smart-io  -w  -ffunction-sections  -fdata-sections  -g  -mdebugger  -Wcast-align  -fno-short-double  -mprocessor=32MX795F512L  -DF_CPU=80000000L  -DARDUINO=23  -D_BOARD_MEGA_    -DMPIDEVER=0x01000305  -DMPIDE=23      -I/Applications/Mpide.app/Contents/Resources/Java/hardware/pic32/cores/pic32   -I/Applications/Mpide.app/Contents/Resources/Java/hardware/pic32/variants/Max32    /Applications/Mpide.app/Contents/Resources/Java/hardware/pic32/cores/pic32/pins_arduino.c  -o  /var/folders/f0/q5xc0_tj4jv4ljd_lfz4pg5h0000gn/T/build2117742867436922308.tmp/pins_arduino.c.o
/Applications/mpide.app/Contents/Resources/Java/hardware/pic32/compiler/pic32-tools/bin/pic32-gcc  -O2  -c  -mno-smart-io  -w  -ffunction-sections  -fdata-sections  -g  -mdebugger  -Wcast-align  -fno-short-double  -mprocessor=32MX795F512L  -DF_CPU=80000000L  -DARDUINO=23  -D_BOARD_MEGA_    -DMPIDEVER=0x01000305  -DMPIDE=23      -I/Applications/Mpide.app/Contents/Resources/Java/hardware/pic32/cores/pic32   -I/Applications/Mpide.app/Contents/Resources/Java/hardware/pic32/variants/Max32    /Applications/Mpide.app/Contents/Resources/Java/hardware/pic32/cores/pic32/task_manager.c  -o  /var/folders/f0/q5xc0_tj4jv4ljd_lfz4pg5h0000gn/T/build2117742867436922308.tmp/task_manager.c.o
/Applications/mpide.app/Contents/Resources/Java/hardware/pic32/compiler/pic32-tools/bin/pic32-gcc  -O2  -c  -mno-smart-io  -w  -ffunction-sections  -fdata-sections  -g  -mdebugger  -Wcast-align  -fno-short-double  -mprocessor=32MX795F512L  -DF_CPU=80000000L  -DARDUINO=23  -D_BOARD_MEGA_    -DMPIDEVER=0x01000305  -DMPIDE=23      -I/Applications/Mpide.app/Contents/Resources/Java/hardware/pic32/cores/pic32   -I/Applications/Mpide.app/Contents/Resources/Java/hardware/pic32/variants/Max32    /Applications/Mpide.app/Contents/Resources/Java/hardware/pic32/cores/pic32/WInterrupts.c  -o  /var/folders/f0/q5xc0_tj4jv4ljd_lfz4pg5h0000gn/T/build2117742867436922308.tmp/WInterrupts.c.o
/Applications/mpide.app/Contents/Resources/Java/hardware/pic32/compiler/pic32-tools/bin/pic32-gcc  -O2  -c  -mno-smart-io  -w  -ffunction-sections  -fdata-sections  -g  -mdebugger  -Wcast-align  -fno-short-double  -mprocessor=32MX795F512L  -DF_CPU=80000000L  -DARDUINO=23  -D_BOARD_MEGA_    -DMPIDEVER=0x01000305  -DMPIDE=23      -I/Applications/Mpide.app/Contents/Resources/Java/hardware/pic32/cores/pic32   -I/Applications/Mpide.app/Contents/Resources/Java/hardware/pic32/variants/Max32    /Applications/Mpide.app/Contents/Resources/Java/hardware/pic32/cores/pic32/wiring.c  -o  /var/folders/f0/q5xc0_tj4jv4ljd_lfz4pg5h0000gn/T/build2117742867436922308.tmp/wiring.c.o
/Applications/mpide.app/Contents/Resources/Java/hardware/pic32/compiler/pic32-tools/bin/pic32-gcc  -O2  -c  -mno-smart-io  -w  -ffunction-sections  -fdata-sections  -g  -mdebugger  -Wcast-align  -fno-short-double  -mprocessor=32MX795F512L  -DF_CPU=80000000L  -DARDUINO=23  -D_BOARD_MEGA_    -DMPIDEVER=0x01000305  -DMPIDE=23      -I/Applications/Mpide.app/Contents/Resources/Java/hardware/pic32/cores/pic32   -I/Applications/Mpide.app/Contents/Resources/Java/hardware/pic32/variants/Max32    /Applications/Mpide.app/Contents/Resources/Java/hardware/pic32/cores/pic32/wiring_analog.c  -o  /var/folders/f0/q5xc0_tj4jv4ljd_lfz4pg5h0000gn/T/build2117742867436922308.tmp/wiring_analog.c.o
/Applications/mpide.app/Contents/Resources/Java/hardware/pic32/compiler/pic32-tools/bin/pic32-gcc  -O2  -c  -mno-smart-io  -w  -ffunction-sections  -fdata-sections  -g  -mdebugger  -Wcast-align  -fno-short-double  -mprocessor=32MX795F512L  -DF_CPU=80000000L  -DARDUINO=23  -D_BOARD_MEGA_    -DMPIDEVER=0x01000305  -DMPIDE=23      -I/Applications/Mpide.app/Contents/Resources/Java/hardware/pic32/cores/pic32   -I/Applications/Mpide.app/Contents/Resources/Java/hardware/pic32/variants/Max32    /Applications/Mpide.app/Contents/Resources/Java/hardware/pic32/cores/pic32/wiring_digital.c  -o  /var/folders/f0/q5xc0_tj4jv4ljd_lfz4pg5h0000gn/T/build2117742867436922308.tmp/wiring_digital.c.o
/Applications/mpide.app/Contents/Resources/Java/hardware/pic32/compiler/pic32-tools/bin/pic32-gcc  -O2  -c  -mno-smart-io  -w  -ffunction-sections  -fdata-sections  -g  -mdebugger  -Wcast-align  -fno-short-double  -mprocessor=32MX795F512L  -DF_CPU=80000000L  -DARDUINO=23  -D_BOARD_MEGA_    -DMPIDEVER=0x01000305  -DMPIDE=23      -I/Applications/Mpide.app/Contents/Resources/Java/hardware/pic32/cores/pic32   -I/Applications/Mpide.app/Contents/Resources/Java/hardware/pic32/variants/Max32    /Applications/Mpide.app/Contents/Resources/Java/hardware/pic32/cores/pic32/wiring_pulse.c  -o  /var/folders/f0/q5xc0_tj4jv4ljd_lfz4pg5h0000gn/T/build2117742867436922308.tmp/wiring_pulse.c.o
/Applications/mpide.app/Contents/Resources/Java/hardware/pic32/compiler/pic32-tools/bin/pic32-gcc  -O2  -c  -mno-smart-io  -w  -ffunction-sections  -fdata-sections  -g  -mdebugger  -Wcast-align  -fno-short-double  -mprocessor=32MX795F512L  -DF_CPU=80000000L  -DARDUINO=23  -D_BOARD_MEGA_    -DMPIDEVER=0x01000305  -DMPIDE=23      -I/Applications/Mpide.app/Contents/Resources/Java/hardware/pic32/cores/pic32   -I/Applications/Mpide.app/Contents/Resources/Java/hardware/pic32/variants/Max32    /Applications/Mpide.app/Contents/Resources/Java/hardware/pic32/cores/pic32/wiring_shift.c  -o  /var/folders/f0/q5xc0_tj4jv4ljd_lfz4pg5h0000gn/T/build2117742867436922308.tmp/wiring_shift.c.o
/Applications/mpide.app/Contents/Resources/Java/hardware/pic32/compiler/pic32-tools/bin/pic32-gcc  -O2  -c  -mno-smart-io  -w  -ffunction-sections  -fdata-sections  -g  -mdebugger  -Wcast-align  -fno-short-double  -mprocessor=32MX795F512L  -DF_CPU=80000000L  -DARDUINO=23  -D_BOARD_MEGA_    -DMPIDEVER=0x01000305  -DMPIDE=23      -I/Applications/Mpide.app/Contents/Resources/Java/hardware/pic32/cores/pic32   -I/Applications/Mpide.app/Contents/Resources/Java/hardware/pic32/variants/Max32    /Applications/Mpide.app/Contents/Resources/Java/hardware/pic32/cores/pic32/WSystem.c  -o  /var/folders/f0/q5xc0_tj4jv4ljd_lfz4pg5h0000gn/T/build2117742867436922308.tmp/WSystem.c.o
/Applications/mpide.app/Contents/Resources/Java/hardware/pic32/compiler/pic32-tools/bin/pic32-g++  -O2  -c  -mno-smart-io  -w  -fno-exceptions  -ffunction-sections  -fdata-sections  -g  -mdebugger  -Wcast-align  -fno-short-double  -mprocessor=32MX795F512L  -DF_CPU=80000000L  -DARDUINO=23  -D_BOARD_MEGA_    -DMPIDEVER=0x01000305  -DMPIDE=23      -I/Applications/Mpide.app/Contents/Resources/Java/hardware/pic32/cores/pic32   -I/Applications/Mpide.app/Contents/Resources/Java/hardware/pic32/variants/Max32    /Applications/Mpide.app/Contents/Resources/Java/hardware/pic32/cores/pic32/HardwareSerial.cpp  -o  /var/folders/f0/q5xc0_tj4jv4ljd_lfz4pg5h0000gn/T/build2117742867436922308.tmp/HardwareSerial.cpp.o
/Applications/mpide.app/Contents/Resources/Java/hardware/pic32/compiler/pic32-tools/bin/pic32-g++  -O2  -c  -mno-smart-io  -w  -fno-exceptions  -ffunction-sections  -fdata-sections  -g  -mdebugger  -Wcast-align  -fno-short-double  -mprocessor=32MX795F512L  -DF_CPU=80000000L  -DARDUINO=23  -D_BOARD_MEGA_    -DMPIDEVER=0x01000305  -DMPIDE=23      -I/Applications/Mpide.app/Contents/Resources/Java/hardware/pic32/cores/pic32   -I/Applications/Mpide.app/Contents/Resources/Java/hardware/pic32/variants/Max32    /Applications/Mpide.app/Contents/Resources/Java/hardware/pic32/cores/pic32/main.cpp  -o  /var/folders/f0/q5xc0_tj4jv4ljd_lfz4pg5h0000gn/T/build2117742867436922308.tmp/main.cpp.o
/Applications/mpide.app/Contents/Resources/Java/hardware/pic32/compiler/pic32-tools/bin/pic32-g++  -O2  -c  -mno-smart-io  -w  -fno-exceptions  -ffunction-sections  -fdata-sections  -g  -mdebugger  -Wcast-align  -fno-short-double  -mprocessor=32MX795F512L  -DF_CPU=80000000L  -DARDUINO=23  -D_BOARD_MEGA_    -DMPIDEVER=0x01000305  -DMPIDE=23      -I/Applications/Mpide.app/Contents/Resources/Java/hardware/pic32/cores/pic32   -I/Applications/Mpide.app/Contents/Resources/Java/hardware/pic32/variants/Max32    /Applications/Mpide.app/Contents/Resources/Java/hardware/pic32/cores/pic32/Print.cpp  -o  /var/folders/f0/q5xc0_tj4jv4ljd_lfz4pg5h0000gn/T/build2117742867436922308.tmp/Print.cpp.o
/Applications/mpide.app/Contents/Resources/Java/hardware/pic32/compiler/pic32-tools/bin/pic32-g++  -O2  -c  -mno-smart-io  -w  -fno-exceptions  -ffunction-sections  -fdata-sections  -g  -mdebugger  -Wcast-align  -fno-short-double  -mprocessor=32MX795F512L  -DF_CPU=80000000L  -DARDUINO=23  -D_BOARD_MEGA_    -DMPIDEVER=0x01000305  -DMPIDE=23      -I/Applications/Mpide.app/Contents/Resources/Java/hardware/pic32/cores/pic32   -I/Applications/Mpide.app/Contents/Resources/Java/hardware/pic32/variants/Max32    /Applications/Mpide.app/Contents/Resources/Java/hardware/pic32/cores/pic32/Tone.cpp  -o  /var/folders/f0/q5xc0_tj4jv4ljd_lfz4pg5h0000gn/T/build2117742867436922308.tmp/Tone.cpp.o
/Applications/mpide.app/Contents/Resources/Java/hardware/pic32/compiler/pic32-tools/bin/pic32-g++  -O2  -c  -mno-smart-io  -w  -fno-exceptions  -ffunction-sections  -fdata-sections  -g  -mdebugger  -Wcast-align  -fno-short-double  -mprocessor=32MX795F512L  -DF_CPU=80000000L  -DARDUINO=23  -D_BOARD_MEGA_    -DMPIDEVER=0x01000305  -DMPIDE=23      -I/Applications/Mpide.app/Contents/Resources/Java/hardware/pic32/cores/pic32   -I/Applications/Mpide.app/Contents/Resources/Java/hardware/pic32/variants/Max32    /Applications/Mpide.app/Contents/Resources/Java/hardware/pic32/cores/pic32/WMath.cpp  -o  /var/folders/f0/q5xc0_tj4jv4ljd_lfz4pg5h0000gn/T/build2117742867436922308.tmp/WMath.cpp.o
/Applications/mpide.app/Contents/Resources/Java/hardware/pic32/compiler/pic32-tools/bin/pic32-g++  -O2  -c  -mno-smart-io  -w  -fno-exceptions  -ffunction-sections  -fdata-sections  -g  -mdebugger  -Wcast-align  -fno-short-double  -mprocessor=32MX795F512L  -DF_CPU=80000000L  -DARDUINO=23  -D_BOARD_MEGA_    -DMPIDEVER=0x01000305  -DMPIDE=23      -I/Applications/Mpide.app/Contents/Resources/Java/hardware/pic32/cores/pic32   -I/Applications/Mpide.app/Contents/Resources/Java/hardware/pic32/variants/Max32    /Applications/Mpide.app/Contents/Resources/Java/hardware/pic32/cores/pic32/WString.cpp  -o  /var/folders/f0/q5xc0_tj4jv4ljd_lfz4pg5h0000gn/T/build2117742867436922308.tmp/WString.cpp.o
/Applications/mpide.app/Contents/Resources/Java/hardware/pic32/compiler/pic32-tools/bin/pic32-ar  rcs  /var/folders/f0/q5xc0_tj4jv4ljd_lfz4pg5h0000gn/T/build2117742867436922308.tmp/core.a  /var/folders/f0/q5xc0_tj4jv4ljd_lfz4pg5h0000gn/T/build2117742867436922308.tmp/cpp-startup.S.o
/Applications/mpide.app/Contents/Resources/Java/hardware/pic32/compiler/pic32-tools/bin/pic32-ar  rcs  /var/folders/f0/q5xc0_tj4jv4ljd_lfz4pg5h0000gn/T/build2117742867436922308.tmp/core.a  /var/folders/f0/q5xc0_tj4jv4ljd_lfz4pg5h0000gn/T/build2117742867436922308.tmp/crti.S.o
/Applications/mpide.app/Contents/Resources/Java/hardware/pic32/compiler/pic32-tools/bin/pic32-ar  rcs  /var/folders/f0/q5xc0_tj4jv4ljd_lfz4pg5h0000gn/T/build2117742867436922308.tmp/core.a  /var/folders/f0/q5xc0_tj4jv4ljd_lfz4pg5h0000gn/T/build2117742867436922308.tmp/crtn.S.o
/Applications/mpide.app/Contents/Resources/Java/hardware/pic32/compiler/pic32-tools/bin/pic32-ar  rcs  /var/folders/f0/q5xc0_tj4jv4ljd_lfz4pg5h0000gn/T/build2117742867436922308.tmp/core.a  /var/folders/f0/q5xc0_tj4jv4ljd_lfz4pg5h0000gn/T/build2117742867436922308.tmp/exceptions.c.o
/Applications/mpide.app/Contents/Resources/Java/hardware/pic32/compiler/pic32-tools/bin/pic32-ar  rcs  /var/folders/f0/q5xc0_tj4jv4ljd_lfz4pg5h0000gn/T/build2117742867436922308.tmp/core.a  /var/folders/f0/q5xc0_tj4jv4ljd_lfz4pg5h0000gn/T/build2117742867436922308.tmp/HardwareSerial_cdcacm.c.o
/Applications/mpide.app/Contents/Resources/Java/hardware/pic32/compiler/pic32-tools/bin/pic32-ar  rcs  /var/folders/f0/q5xc0_tj4jv4ljd_lfz4pg5h0000gn/T/build2117742867436922308.tmp/core.a  /var/folders/f0/q5xc0_tj4jv4ljd_lfz4pg5h0000gn/T/build2117742867436922308.tmp/HardwareSerial_usb.c.o
/Applications/mpide.app/Contents/Resources/Java/hardware/pic32/compiler/pic32-tools/bin/pic32-ar  rcs  /var/folders/f0/q5xc0_tj4jv4ljd_lfz4pg5h0000gn/T/build2117742867436922308.tmp/core.a  /var/folders/f0/q5xc0_tj4jv4ljd_lfz4pg5h0000gn/T/build2117742867436922308.tmp/pins_arduino.c.o
/Applications/mpide.app/Contents/Resources/Java/hardware/pic32/compiler/pic32-tools/bin/pic32-ar  rcs  /var/folders/f0/q5xc0_tj4jv4ljd_lfz4pg5h0000gn/T/build2117742867436922308.tmp/core.a  /var/folders/f0/q5xc0_tj4jv4ljd_lfz4pg5h0000gn/T/build2117742867436922308.tmp/task_manager.c.o
/Applications/mpide.app/Contents/Resources/Java/hardware/pic32/compiler/pic32-tools/bin/pic32-ar  rcs  /var/folders/f0/q5xc0_tj4jv4ljd_lfz4pg5h0000gn/T/build2117742867436922308.tmp/core.a  /var/folders/f0/q5xc0_tj4jv4ljd_lfz4pg5h0000gn/T/build2117742867436922308.tmp/WInterrupts.c.o
/Applications/mpide.app/Contents/Resources/Java/hardware/pic32/compiler/pic32-tools/bin/pic32-ar  rcs  /var/folders/f0/q5xc0_tj4jv4ljd_lfz4pg5h0000gn/T/build2117742867436922308.tmp/core.a  /var/folders/f0/q5xc0_tj4jv4ljd_lfz4pg5h0000gn/T/build2117742867436922308.tmp/wiring.c.o
/Applications/mpide.app/Contents/Resources/Java/hardware/pic32/compiler/pic32-tools/bin/pic32-ar  rcs  /var/folders/f0/q5xc0_tj4jv4ljd_lfz4pg5h0000gn/T/build2117742867436922308.tmp/core.a  /var/folders/f0/q5xc0_tj4jv4ljd_lfz4pg5h0000gn/T/build2117742867436922308.tmp/wiring_analog.c.o
/Applications/mpide.app/Contents/Resources/Java/hardware/pic32/compiler/pic32-tools/bin/pic32-ar  rcs  /var/folders/f0/q5xc0_tj4jv4ljd_lfz4pg5h0000gn/T/build2117742867436922308.tmp/core.a  /var/folders/f0/q5xc0_tj4jv4ljd_lfz4pg5h0000gn/T/build2117742867436922308.tmp/wiring_digital.c.o
/Applications/mpide.app/Contents/Resources/Java/hardware/pic32/compiler/pic32-tools/bin/pic32-ar  rcs  /var/folders/f0/q5xc0_tj4jv4ljd_lfz4pg5h0000gn/T/build2117742867436922308.tmp/core.a  /var/folders/f0/q5xc0_tj4jv4ljd_lfz4pg5h0000gn/T/build2117742867436922308.tmp/wiring_pulse.c.o
/Applications/mpide.app/Contents/Resources/Java/hardware/pic32/compiler/pic32-tools/bin/pic32-ar  rcs  /var/folders/f0/q5xc0_tj4jv4ljd_lfz4pg5h0000gn/T/build2117742867436922308.tmp/core.a  /var/folders/f0/q5xc0_tj4jv4ljd_lfz4pg5h0000gn/T/build2117742867436922308.tmp/wiring_shift.c.o
/Applications/mpide.app/Contents/Resources/Java/hardware/pic32/compiler/pic32-tools/bin/pic32-ar  rcs  /var/folders/f0/q5xc0_tj4jv4ljd_lfz4pg5h0000gn/T/build2117742867436922308.tmp/core.a  /var/folders/f0/q5xc0_tj4jv4ljd_lfz4pg5h0000gn/T/build2117742867436922308.tmp/WSystem.c.o
/Applications/mpide.app/Contents/Resources/Java/hardware/pic32/compiler/pic32-tools/bin/pic32-ar  rcs  /var/folders/f0/q5xc0_tj4jv4ljd_lfz4pg5h0000gn/T/build2117742867436922308.tmp/core.a  /var/folders/f0/q5xc0_tj4jv4ljd_lfz4pg5h0000gn/T/build2117742867436922308.tmp/HardwareSerial.cpp.o
/Applications/mpide.app/Contents/Resources/Java/hardware/pic32/compiler/pic32-tools/bin/pic32-ar  rcs  /var/folders/f0/q5xc0_tj4jv4ljd_lfz4pg5h0000gn/T/build2117742867436922308.tmp/core.a  /var/folders/f0/q5xc0_tj4jv4ljd_lfz4pg5h0000gn/T/build2117742867436922308.tmp/main.cpp.o
/Applications/mpide.app/Contents/Resources/Java/hardware/pic32/compiler/pic32-tools/bin/pic32-ar  rcs  /var/folders/f0/q5xc0_tj4jv4ljd_lfz4pg5h0000gn/T/build2117742867436922308.tmp/core.a  /var/folders/f0/q5xc0_tj4jv4ljd_lfz4pg5h0000gn/T/build2117742867436922308.tmp/Print.cpp.o
/Applications/mpide.app/Contents/Resources/Java/hardware/pic32/compiler/pic32-tools/bin/pic32-ar  rcs  /var/folders/f0/q5xc0_tj4jv4ljd_lfz4pg5h0000gn/T/build2117742867436922308.tmp/core.a  /var/folders/f0/q5xc0_tj4jv4ljd_lfz4pg5h0000gn/T/build2117742867436922308.tmp/Tone.cpp.o
/Applications/mpide.app/Contents/Resources/Java/hardware/pic32/compiler/pic32-tools/bin/pic32-ar  rcs  /var/folders/f0/q5xc0_tj4jv4ljd_lfz4pg5h0000gn/T/build2117742867436922308.tmp/core.a  /var/folders/f0/q5xc0_tj4jv4ljd_lfz4pg5h0000gn/T/build2117742867436922308.tmp/WMath.cpp.o
/Applications/mpide.app/Contents/Resources/Java/hardware/pic32/compiler/pic32-tools/bin/pic32-ar  rcs  /var/folders/f0/q5xc0_tj4jv4ljd_lfz4pg5h0000gn/T/build2117742867436922308.tmp/core.a  /var/folders/f0/q5xc0_tj4jv4ljd_lfz4pg5h0000gn/T/build2117742867436922308.tmp/WString.cpp.o
/Applications/mpide.app/Contents/Resources/Java/hardware/pic32/compiler/pic32-tools/bin/pic32-g++  -Os  -Wl,--gc-sections  -mdebugger  -mprocessor=32MX795F512L  -o  /var/folders/f0/q5xc0_tj4jv4ljd_lfz4pg5h0000gn/T/build2117742867436922308.tmp/Blink.cpp.elf  /var/folders/f0/q5xc0_tj4jv4ljd_lfz4pg5h0000gn/T/build2117742867436922308.tmp/Blink.cpp.o    /var/folders/f0/q5xc0_tj4jv4ljd_lfz4pg5h0000gn/T/build2117742867436922308.tmp/core.a  -L/var/folders/f0/q5xc0_tj4jv4ljd_lfz4pg5h0000gn/T/build2117742867436922308.tmp  -lm  -T  /Applications/Mpide.app/Contents/Resources/Java/hardware/pic32/cores/pic32/chipKIT-application-32MX795F512.ld  -T/Applications/Mpide.app/Contents/Resources/Java/hardware/pic32/cores/pic32/chipKIT-application-COMMON.ld
/Applications/mpide.app/Contents/Resources/Java/hardware/pic32/compiler/pic32-tools/bin/pic32-objcopy  -O  ihex  -j  .eeprom  --set-section-flags=.eeprom=alloc,load  --no-change-warnings  --change-section-lma  .eeprom=0  /var/folders/f0/q5xc0_tj4jv4ljd_lfz4pg5h0000gn/T/build2117742867436922308.tmp/Blink.cpp.elf  /var/folders/f0/q5xc0_tj4jv4ljd_lfz4pg5h0000gn/T/build2117742867436922308.tmp/Blink.cpp.eep
/Applications/mpide.app/Contents/Resources/Java/hardware/pic32/compiler/pic32-tools/bin/pic32-bin2hex  -a  /var/folders/f0/q5xc0_tj4jv4ljd_lfz4pg5h0000gn/T/build2117742867436922308.tmp/Blink.cpp.elf
Binary sketch size: 4088 bytes (of a 520192 byte maximum)
\end{lstlisting}

\end{document}
